\documentclass[12pt]{article}

\usepackage[margin=1in]{geometry}
\usepackage{amsmath,amssymb,amsthm}
\usepackage{booktabs}
\usepackage{hyperref}
\usepackage[T1]{fontenc}
\usepackage{lmodern}

\hypersetup{
  colorlinks=true,
  linkcolor=blue,
  citecolor=blue,
  urlcolor=blue
}

\title{Connections between Stachurski--Sargent (2024) and\\
Kreps--Porteus (1977, 1979)}
\author{}
\date{March 2026}

\begin{document}
\maketitle

\section{Shared Foundational Project: Extending Dynamic Programming beyond Standard MDPs}

All three papers share a fundamental ambition: \textbf{generalizing the theory of dynamic programming beyond the standard Markov decision process (MDP) framework}, where the objective is to maximize the expected discounted sum of rewards.

\begin{itemize}
\item \textbf{Kreps--Porteus (SIAM 1977)} explicitly states in its abstract that the goal is ``extending these results to criteria other than the usual expected sum of rewards, such as expected utility criteria, certain stochastic games, risk sensitive Markov decision processes, and maximin criteria.''

\item \textbf{Kreps--Porteus (Econometrica 1979)} axiomatizes ``temporal von Neumann--Morgenstern utility,'' a non-additive recursive preference criterion, and shows that dynamic programming still works for it.

\item \textbf{Stachurski--Sargent (2024)} builds a framework encompassing ``nonlinear recursive preferences, desire for robustness, function approximation, Monte Carlo sampling and distributional dynamic programs''---a modern descendant of the same generalization program.
\end{itemize}


\section{Order-Theoretic and Lattice-Theoretic Machinery}

The most direct technical lineage runs through the use of \textbf{partial orders and monotonicity (isotonicity)} as the engine for optimality results.

\begin{itemize}
\item \textbf{Kreps--Porteus (1977)} works with structured value function spaces $V_t$ and $V_t^*$, and assumes that the single-stage operators $H_t^\delta$ are \textbf{isotone} (order-preserving) and \textbf{Lipschitzian} on these spaces (Assumption~R). Their proofs of optimality of structured policies rest on \textbf{monotone convergence} assumptions (IMC, DMC) and \textbf{closure} assumptions (IC, DC, GDC)---all of which are properties of partially ordered sets under pointwise ordering.

\item \textbf{Stachurski--Sargent} makes this order-theoretic foundation fully explicit and abstract. Their ADP is defined as a pair $(V, \mathcal{T})$ where $V$ is a partially ordered set and each $T_\sigma$ is an \textbf{order-preserving self-map} on $V$. Their key assumptions---\textbf{order stability}, \textbf{order continuity}, \textbf{chain completeness}, and \textbf{countable Dedekind completeness}---are abstract generalizations of exactly the monotone convergence and closure assumptions that Kreps--Porteus imposed concretely.
\end{itemize}

\medskip
\noindent
The following table summarizes the correspondences:

\bigskip
\begin{center}
\begin{tabular}{@{}p{0.45\textwidth}p{0.45\textwidth}@{}}
\toprule
\textbf{Kreps--Porteus (1977)} & \textbf{Stachurski--Sargent (2024)} \\
\midrule
Isotonicity of $H_t^\delta$ & Order-preserving policy operators $T_\sigma$ \\[4pt]
Increasing Closure (IC) & Countably chain complete / chain complete $V$ \\[4pt]
Increasing Monotone Convergence (IMC) & Order continuity of $T_\sigma$ \\[4pt]
Decreasing Closure (DC), General Decreasing Closure (GDC) & Chain completeness / Dedekind completeness \\[4pt]
Decreasing Monotone Convergence (DMC) & Order continuity (symmetric treatment) \\
\bottomrule
\end{tabular}
\end{center}


\section{Positive and Negative Dynamic Programming $\leftrightarrow$ Upward and Downward Stability}

Kreps--Porteus (1977) develops two parallel theories:
\begin{itemize}
\item \textbf{Positive problems} (\S2): Value functions are built up from below; the optimal value is the \textbf{least} solution of the optimality equations exceeding a lower bound; \textbf{unimprovable} strategies are optimal.

\item \textbf{Negative problems} (\S3): Value functions are built down from above; the optimal value is the \textbf{greatest} solution below an upper bound; \textbf{conserving} strategies are optimal.
\end{itemize}

\noindent
Stachurski--Sargent's framework provides a unified abstract treatment through the concepts of \textbf{upward stability} ($v \preceq Sv \Rightarrow v \preceq \bar{v}$) and \textbf{downward stability} ($Sv \preceq v \Rightarrow \bar{v} \preceq v$). Their Theorem~5.1 (if the ADP is downward stable and $T$ has a fixed point, fundamental optimality holds) directly parallels the negative-problem results of Kreps--Porteus, while constructive convergence of VFI from below (Theorems~5.3, 5.5) parallels the positive-problem results.


\section{The Bellman Equation, Greedy Policies, and the Principle of Optimality}

All three papers establish the same trio of fundamental results:
\begin{enumerate}
\item[\textbf{(B1)}] Existence of an optimal value function $v^*$ (greatest element of lifetime values).
\item[\textbf{(B2)}] $v^*$ is the unique solution to the Bellman equation.
\item[\textbf{(B3)}] Bellman's principle of optimality: $\sigma$ is optimal $\iff$ $\sigma$ is $v^*$-greedy.
\end{enumerate}

\noindent
In Kreps--Porteus (1977), these correspond to Theorem~4 parts~(d)--(h) for positive problems and Theorems~5--6 for negative problems. In Kreps--Porteus (1979), Theorem~3 establishes an analogous result in the temporal utility setting. Stachurski--Sargent's Lemma~4.2 proves that (B2) and (B3) are equivalent given (B1), and their main theorems (5.1--5.5) give different sufficient conditions for all three.


\section{Kreps--Porteus Temporal Preferences as a Motivating Application for Abstract DPs}

Kreps--Porteus (1979) introduces preferences where the recursion defining value functions is:
\[
v_t(\pi;\, y_t, x_t)
= E_{\delta_t}\!\bigl[u_t^*\bigl(y_t,\, \tilde{z}_t,\,
      v_{t+1}(\pi;\, (y_t, \tilde{z}_t),\, \tilde{x}_{t+1})\bigr)\bigr].
\]
This is \textbf{not} additively separable---the aggregator $u_t^*$ is a general nonlinear function of the continuation value. This class of preferences (later developed by Epstein--Zin (1989) into the widely-used recursive utility framework) is exactly one of the motivating applications for Stachurski--Sargent. Their Section~7.1 on ``Non-EU Discrete Choice'' and their treatment of risk-sensitive preferences (Section~3.2) handle generalizations of Kreps--Porteus temporal preferences by replacing the standard conditional expectation with a \textbf{certainty equivalent operator} $\mathcal{E}$ that can be nonlinear---encompassing both the Kreps--Porteus temporal utility and risk-sensitive/robustness preferences.


\section{Policy Iteration and Strategy Improvement}

Kreps--Porteus (1977) proves that in negative problems, if strategy $\pi$ is a \textbf{one-step improvement} on $\pi'$, then $\pi$ nets a higher value (Theorem~5(h)). This rules out cycling in strategy (policy) iteration. Stachurski--Sargent provide convergence results for \textbf{Howard policy iteration (HPI)}---the modern name for the same procedure---under their abstract conditions (Theorems~5.3--5.5, with finite-step convergence in Theorem~5.4).


\section{The Role of the Porteus (1975) Operator Framework}

Both Kreps--Porteus papers build on Porteus (1975), ``On the Optimality of Structured Policies in Countable Stage Decision Processes,'' which introduced the operator-based approach to dynamic programming with structured policies. The 1977 paper is explicitly ``Part~II'' of this line. Stachurski--Sargent's framework can be seen as a further abstraction of this same operator approach---where Porteus's single-stage operators $H_t^\delta$ evolve into the family of policy operators $\{T_\sigma\}$, and the structured value function spaces $V_t^*$ evolve into the abstract partially ordered set~$V$.


\section{Summary}

The Stachurski--Sargent paper is a modern, abstract descendant of the research program initiated by Kreps and Porteus in the 1970s. Kreps--Porteus pioneered:
\begin{enumerate}
\item[(a)] extending dynamic programming optimality theory beyond expected-sum-of-rewards criteria using order-theoretic (isotonicity/monotone-convergence) methods, and
\item[(b)] showing that non-additive recursive preferences (temporal utility) are amenable to dynamic programming.
\end{enumerate}

\noindent
Stachurski--Sargent takes these ideas to their logical conclusion by working in a fully abstract partially ordered set, unifying positive/negative DP into a single framework via order stability, and encompassing not just Kreps--Porteus preferences but also distributional DP, empirical DP, function approximation, and other modern applications.


\bigskip
\subsection*{Papers discussed}

\begin{enumerate}
\item D.\,M.\ Kreps and E.\,L.\ Porteus, ``On the Optimality of Structured Policies in Countable Stage Decision Processes.\ II: Positive and Negative Problems,'' \textit{SIAM Journal on Applied Mathematics}, 32(2):457--466, 1977.

\item D.\,M.\ Kreps and E.\,L.\ Porteus, ``Dynamic Choice Theory and Dynamic Programming,'' \textit{Econometrica}, 47(1):91--100, 1979.

\item J.\ Stachurski and T.\,J.\ Sargent, ``Dynamic Programs on Partially Ordered Sets,'' 2024. Available at \url{http://www.tomsargent.com/research/adps_3.pdf}.
\end{enumerate}

\end{document}
